%% Chapter 3, Section 3.1: Classical Monte Carlo
%% A First Course in Monte Carlo Methods — Sanz-Alonso & Al-Ghattas

\section{Classical Monte Carlo}
\label{sec:3.1}

Classical Monte Carlo integration presupposes that we can obtain $N$ independent
draws from the target distribution. Then, the integral of interest $\calI_f[h]$
is approximated by the average of the values of $h$ at those draws. We refer to
the method, summarized in Algorithm~\ref{alg:3.1} below, as \emph{classical}
Monte Carlo to distinguish it from other Monte Carlo methods. In the literature
it is often simply called Monte Carlo, standard Monte Carlo, or vanilla Monte
Carlo.

\begin{algorithm}[H]
\caption{Classical Monte Carlo Integration}
\label{alg:3.1}
\begin{algorithmic}[1]
\Require Target distribution $f$, test function $h$, sample size $N$.
\State Sample $X^{(1)}, \ldots, X^{(N)} \iid f$.
\Ensure Monte Carlo estimator
  $\calI_f^{\mathrm{mc}}[h] := \dfrac{1}{N} \displaystyle\sum_{n=1}^{N}
   h\!\left(X^{(n)}\right) \approx \calI_f[h]$.
\end{algorithmic}
\end{algorithm}

The following result shows that the classical Monte Carlo estimator is unbiased
and that its mean squared error, which agrees with its variance, is of order
$1/N$.

\begin{theorem}[Classical Monte Carlo Error]
\label{thm:3.1}
The following holds:
\begin{enumerate}
  \item $\Expect\!\left[\calI_f^{\mathrm{mc}}[h]\right] = \calI_f[h]$;
  \item $\Var\!\left[\calI_f^{\mathrm{mc}}[h]\right]
         = \dfrac{1}{N}\,\Var_{X \sim f}\!\left[h(X)\right]$.
\end{enumerate}
\end{theorem}

\begin{proof}
For the first claim, linearity of expectation and the fact that $X^{(n)} \sim f$
give
\[
  \Expect\!\left[\calI_f^{\mathrm{mc}}[h]\right]
  = \Expect\!\left[\frac{1}{N}\sum_{n=1}^{N} h\!\left(X^{(n)}\right)\right]
  = \frac{1}{N}\sum_{n=1}^{N} \Expect\!\left[h\!\left(X^{(n)}\right)\right]
  = \Expect_{X \sim f}\!\left[h(X)\right]
  = \calI_f[h].
\]
For the second claim, note that since $X^{(n)} \iid f$,
\[
  \Var\!\left[\calI_f^{\mathrm{mc}}[h]\right]
  = \frac{1}{N^2}\sum_{n=1}^{N} \Var\!\left[h\!\left(X^{(n)}\right)\right]
  = \frac{1}{N^2}\, N\, \Var_{X \sim f}\!\left[h(X)\right]
  = \frac{1}{N}\,\Var_{X \sim f}\!\left[h(X)\right]. \qedhere
\]
\end{proof}

Theorem~\ref{thm:3.1} can be used to construct confidence intervals for
$\calI_f^{\mathrm{mc}}[h]$ as an estimator of $\calI_f[h]$. To do so, note
that by the central limit theorem, as $N \to \infty$,
\begin{equation}
\label{eq:3.2}
  \sqrt{N}\,\frac{\calI_f^{\mathrm{mc}}[h] - \calI_f[h]}%
               {\sqrt{\Var_{X \sim f}[h(X)]}}
  \cvd \Normal(0,1).
\end{equation}
For $\alpha \in (0,1)$ define the $z$-score $z_{\alpha/2}$ by the requirement
that $\Prob(-z_{\alpha/2} < Z < z_{\alpha/2}) = 1 - \alpha$, where
$Z \sim \Normal(0,1)$. Then, it follows from \eqref{eq:3.2} that
\[
  \left(
    \calI_f^{\mathrm{mc}}[h] - z_{\alpha/2}\sqrt{\frac{\Var_{X \sim f}[h(X)]}{N}},\;
    \calI_f^{\mathrm{mc}}[h] + z_{\alpha/2}\sqrt{\frac{\Var_{X \sim f}[h(X)]}{N}}
  \right)
\]
is an approximate $1-\alpha$ confidence interval for $\calI_f[h]$. The variance
term can also be approximated by Monte Carlo:
\begin{equation}
\label{eq:3.3}
  \Var_{X \sim f}\!\left[h(X)\right]
  = \Expect\!\left[\bigl(h(X) - \Expect_{X \sim f}[h(X)]\bigr)^2\right]
  \approx
  \frac{1}{N}\sum_{n=1}^{N}
    \bigl(h(X^{(n)}) - \calI_f^{\mathrm{mc}}[h]\bigr)^2.
\end{equation}

The sample size $N$ is typically large in Monte Carlo simulations, and hence
large $N$ asymptotics, needed to invoke the central limit theorem \eqref{eq:3.2}
and to ensure the convergence of the variance approximation \eqref{eq:3.3}, are
well justified.

\medskip
\noindent\textbf{Pros and Cons.}
The standard error of the classical Monte Carlo estimator is of order
$N^{-1/2}$, which implies a very slow rate of convergence compared to most
deterministic quadrature methods: increasing by a factor of $100$ the number of
samples only reduces the expected error by a factor of $10$. Another drawback of
Monte Carlo methods is that the error can only be understood probabilistically,
while for deterministic quadrature rules deterministic bounds can be derived.
On the bright side, Theorem~\ref{thm:3.1} does not make any assumption on the
state space $E$ (or its dimension) and the only implicit assumption on $h$ is
that $\Expect_{X \sim f}[h(X)^2] < \infty$. In contrast, classical error
bounds for deterministic quadrature rules in $\R^d$ need the number $N$ of
function evaluations to scale as $N^d$ as $d \to \infty$. Moreover,
differentiability conditions on $h$ often needed for deterministic methods play
no role in Theorem~\ref{thm:3.1}. The insensitivity of classical Monte Carlo to
dimension and roughness of $f$ and $h$ makes it appealing in cases where
deterministic methods struggle. It is important to note, however, that the
classical Monte Carlo method presupposes that sampling from $f$ (without
sampling error) is possible. When this is not true (which is often the case in
applications) the error of Monte Carlo methods that rely on approximate target
samples may be sensitive to the dimension of the state space.
